% =========================================================
% Strong Induction
% =========================================================
\begin{tcolorbox}[colback=propbox, colframe=propborder, arc=2pt,
  left=6pt, right=6pt, top=4pt, bottom=4pt,
  title={\small\textbf{Theorem: Strong Induction}},
  fonttitle=\small\bfseries]
Let $P(n)$ be a statement depending on $n \in \mathbb{N}$.
Suppose that for every $n \in \mathbb{N}$:
\[
\Big(\forall k < n,\; P(k)\Big) \Rightarrow P(n).
\]
Then $P(n)$ holds for all $n \in \mathbb{N}$.
\end{tcolorbox}

\begin{remark}[Strong and weak induction prove the same statements]
Strong induction is sometimes described as ``more powerful'' than weak
induction, but this is imprecise. Both principles are logically
equivalent (they both follow from P5 and each implies the other).
Strong induction is not more powerful in the sense of proving statements
that weak induction cannot --- every statement provable by one is
provable by the other.

What strong induction offers is convenience: when the proof of $P(n+1)$
genuinely requires knowing $P(k)$ for some $k < n$ that is not $n-1$,
strong induction supplies that information directly. Rephrasing the proof
as weak induction would require artificially strengthening $P(n)$ to
``$P(k)$ for all $k \leq n$'', which is precisely the hypothesis that
strong induction builds in automatically.

Use strong induction when the step needs information from more than one
earlier case. Use weak induction when $P(n)$ alone suffices.
\end{remark}

\begin{remark}[The base case in strong induction]
When $n = 0$, the hypothesis $\forall k < 0,\; P(k)$ is vacuously
true: there are no natural numbers less than $0$. So the step for
$n = 0$ requires proving $P(0)$ with no inductive hypothesis at all
--- this is the base case, and it must be proved directly.

In practice, always write the base case explicitly, even though it is
implicit in the universal statement. Students who skip the base case in
strong induction often write ``by the inductive hypothesis applied to
$n-1$'' when $n = 0$, which is an error: $n - 1 = -1$ is not a natural
number and has no inductive hypothesis.
\end{remark}

\begin{remark}[When to use strong induction: the diagnostic]
The diagnostic for strong induction is simple. Write out the inductive
step: $P(n+1) = f(P(?), P(?), \ldots)$, and ask what earlier values
of $P$ you need. If the answer includes anything other than exactly $P(n)$,
use strong induction.

Canonical situations where strong induction is needed:
\begin{itemize}[itemsep=2pt]
  \item Recursive sequences with two prior terms: $a_{n+1} = a_n + a_{n-1}$
    requires both $P(n)$ and $P(n-1)$.
  \item Prime factorization: writing $n = ab$ with $a, b < n$
    requires $P(a)$ and $P(b)$ where $a$ and $b$ are arbitrary.
  \item Any argument that splits $n+1$ into parts: you need $P$ at the
    part values, which are not predetermined.
\end{itemize}
\end{remark}

\begin{tcolorbox}[colback=gray!4, colframe=gray!35, arc=2pt,
  left=6pt, right=6pt, top=4pt, bottom=4pt,
  title={\small\textbf{Template: Strong Induction Proof}},
  fonttitle=\small\bfseries]
Let $P(n)$ denote: [\emph{write it out}].

\medskip
\noindent\textbf{Base case} ($n = [\text{start}]$).
[\emph{Verify directly, with no inductive hypothesis.}]

\medskip
\noindent\textbf{Inductive step.}
Let $n > [\text{start}]$ and assume $P(k)$ holds for all
$[\text{start}] \leq k < n$.
\textit{(Strong inductive hypothesis.)}
We show $P(n)$ holds.

[\emph{Use $P(k)$ for some specific $k < n$.}]

Hence $P(n)$ holds.

By strong induction, $P(n)$ holds for all $n \geq [\text{start}]$. \qed
\end{tcolorbox}

\begin{example}[Every integer $n \geq 2$ has a prime factor]
\textit{Proof.}
Let $P(n)$: ``$n$ has a prime factor.''

\medskip
\noindent\textbf{Base case} ($n = 2$).
$2$ is prime, so $2$ is a prime factor of itself. $P(2)$ holds.

\medskip
\noindent\textbf{Inductive step.}
Let $n > 2$ and assume $P(k)$ holds for all $2 \leq k < n$.

\textbf{Case 1:} $n$ is prime. Then $n$ is its own prime factor,
and $P(n)$ holds.

\textbf{Case 2:} $n$ is composite. Write $n = ab$ with $2 \leq a, b < n$.
By the strong inductive hypothesis applied to $a$, $a$ has a prime
factor $p$. Then $p \mid a$ and $a \mid n$, so $p \mid n$.
Hence $P(n)$ holds.

In both cases $P(n)$ holds. By strong induction, $P(n)$ holds for
all $n \geq 2$. \qed

\medskip
\noindent
Why weak induction cannot work directly: in Case~2, the value $a$
satisfies $2 \leq a < n$ but is otherwise arbitrary. It is not
necessarily $n-1$. Weak induction supplies $P(n-1)$, which is useless
when $a \neq n-1$. Strong induction supplies $P(k)$ for all $k < n$,
covering whatever value $a$ takes.
\end{example}
